\chapter{Bussen}
\label{ch:buss}

{\large\itshape Två ledningar, många noder, och en asymmetri som avgör allt.\par}

\begin{chapterbox}{I det här kapitlet}
\begin{itemize}
  \item Varför punkt-till-punkt slutar fungera när ett system växer.
  \item Den delade bussen, och vad den kostar.
  \item Dominanta och recessiva bitar, och den logiska OCH-funktionen bussen utför.
  \item Differentiell signalering och terminering, så långt den som bygger en nod behöver det.
\end{itemize}
\end{chapterbox}

\section{Problemet}

Ett inbyggt system med fem enheter som alla behöver prata med varandra behöver tio förbindelser
om varje par kopplas ihop direkt. Med tio enheter är det fyrtiofem. Formeln är $n(n-1)/2$, och
den är obeveklig: kablaget växer kvadratiskt medan nyttan växer linjärt.

I en bil är det inte en teoretisk invändning. Motorstyrning, ABS, instrumentpanel, växellåda,
krockkuddar, klimatanläggning --- och varje förbindelse är koppar som väger, kostar, tar plats och
kan gå av.

\textbf{Lösningen är att sluta koppla ihop enheter, och i stället koppla in dem.} En enda buss,
som alla noder delar:

\bookfigure{bus_topology}{Punkt-till-punkt mot delad buss. Fem noder kräver tio förbindelser i
det vänstra fallet och en buss i det högra.}{fig:topology}

Det som var $n(n-1)/2$ förbindelser blir $n$ anslutningar till en gemensam ledning. Priset är
att bussen nu är en delad resurs, och delade resurser måste fördelas.

\section{Vad en delad buss kostar}

Tre problem uppstår i samma ögonblick som två noder delar en ledning, och CAN:s hela konstruktion
är svaret på dem:

\begin{booktable}{@{}lL@{}}
\thead{Problem} & \thead{CAN:s svar} \\ \midrule
Alla hör allt, så en mottagare måste veta vad som angår den & Identifieraren beskriver
  \emph{innehållet}, och varje nod filtrerar själv (\kapref{ch:frameformat}) \\
Två noder kan börja sända samtidigt & Icke-förstörande arbitrering (\kapref{ch:arbitrering}) \\
En störning på ledningen träffar alla & Bitstoppning och CRC-15 (\kapref{ch:stuffing}), plus
  felhantering (\kapref{ch:fel}) \\
\end{booktable}

Lägg märke till vad som \emph{inte} står i tabellen: adressering. CAN har ingen
destinationsadress. Det är inte en förenkling utan ett designval, och \kapref{ch:frameformat} går igenom
varför.

\section{Dominant och recessiv}

Här ligger den enda idé man måste ha helt klar för sig för att resten av CAN ska bli begripligt.

En CAN-buss har två tillstånd, och de är \textbf{inte} symmetriska:

\begin{narrowtable}{@{}lll@{}}
\thead{Tillstånd} & \thead{Logisk nivå} & \thead{Vinner?} \\ \midrule
\dom{} & 0 & Ja \\
\rec{} & 1 & Nej \\
\end{narrowtable}

Regeln är enkel och får konsekvenser hela boken igenom:

\begin{aside}
\textbf{Bussen är dominant om \emph{minst en} nod driver den dominant. Den är recessiv bara om
\emph{alla} noder släpper den.}
\end{aside}

Elektriskt är det en öppen dränering med ett uppdragsmotstånd: en nod kan bara \emph{dra ned}
ledningen, aldrig driva upp den. Släpper alla, dras den upp av sig själv.

Logiskt är bussen alltså en OCH-grind över alla noders utgångar:

\bookfigure{wired_and}{Bussen som wired-AND. Tre noder, och bussens nivå som funktion av vad de
driver: en enda dominant nod räcker för att bussen ska bli dominant.}{fig:wiredand}

\subsection{Varför asymmetrin behövs}

Anta i stället en symmetrisk buss, där en etta och en nolla som möts ger ett odefinierat
mellanläge. Då är en kollision en \emph{förstörelse}: båda framesen blir obrukbara, båda noderna
måste upptäcka det, vänta olika länge och försöka igen. Det är så Ethernet gjorde det, och det
är därför Ethernet behövde slumpmässig backoff.

Med asymmetrin blir en kollision i stället en \emph{jämförelse}. Noden som sänder recessivt och
läser tillbaka dominant vet omedelbart två saker: att någon annan också sänder, och att den
andra har företräde. Den drar sig ur, tyst, mitt i framen --- och den vinnande nodens frame fortsätter
utan ett enda förlorat bitintervall.

Det kallas \textbf{icke-förstörande arbitrering}, och det är CAN:s mest eleganta egenskap.
\Kapref{ch:arbitrering} arbetar igenom den bit för bit.

\subsection{Att läsa tillbaka det man sänder}

En förutsättning som är lätt att gå förbi: \textbf{varje sändande nod läser hela tiden av bussen
medan den sänder.} Det är så den upptäcker att den förlorat arbitreringen, och det är också så
den upptäcker att bussen inte gör som den blev tillsagd.

Den förenklade kontrollern gör detta under arbitreringsfältet. En fullständig CAN-kontroller gör
det under \emph{hela} framen, och rapporterar ett \emph{bitfel} så fort det den läser tillbaka
skiljer sig från det den sände. Skillnaden står i \kapref{ch:kursen}.

\section{Ledningarna}

CAN är \textbf{differentiellt}: signalen ligger som en spänningsskillnad mellan två ledare,
\code{CAN_H} och \code{CAN_L}, tvinnade om varandra.

\begin{narrowtable}{@{}lll@{}}
\thead{Tillstånd} & \thead{\code{CAN_H}} & \thead{\code{CAN_L}} \\ \midrule
\rec{} & 2,5 V & 2,5 V \\
\dom{} & 3,5 V & 1,5 V \\
\end{narrowtable}

I recessivt läge driver ingen, och båda ledarna vilar på samma nivå --- skillnaden är noll. I
dominant läge driver den sändande noden isär dem till ungefär två volt.

Poängen med differentiell signalering är störimmunitet. En störning som slår mot kabeln träffar
båda ledarna ungefär lika mycket, och eftersom mottagaren bara mäter \emph{skillnaden} tar den
inte notis om den. Att ledarna är tvinnade är vad som gör att störningen verkligen träffar båda
lika.

\subsection{Terminering}

En CAN-buss termineras med \textbf{120~\textohm{} i vardera änden}, inte i mitten, och inte fler än
två motstånd oavsett hur många noder som sitter på bussen.

Motstånden matchar kabelns karakteristiska impedans så att en signal som når änden absorberas i
stället för att studsa tillbaka. En felterminerad buss fungerar ofta ändå på en kort kabel vid
låg hastighet, och slutar fungera precis när man börjat lita på den --- vilket gör felet till ett
av de obehagligare.

\begin{aside}
\note[Vid bänken] Ett snabbt sätt att kontrollera termineringen: mät resistansen över
\code{CAN_H} och \code{CAN_L} med bussen strömlös. Två parallella 120~\textohm{} ska ge ungefär
60~\textohm. Får ni 120 saknas ett motstånd; får ni 40 sitter det tre.
\end{aside}

\section{Hastighet och längd}

CAN:s bithastighet och bussens längd hänger ihop, och skälet är arbitreringen: under
arbitreringsfältet måste en bit hinna ut till bussens bortersta nod \emph{och} tillbaka innan
biten är slut, annars kan inte noderna vara överens om vem som vann.

\begin{narrowtable}{@{}ll@{}}
\thead{Bithastighet} & \thead{Ungefärlig maximal längd} \\ \midrule
1 Mbit/s & 40 m \\
500 kbit/s & 100 m \\
125 kbit/s & 500 m \\
\end{narrowtable}

Den förenklade kontrollern kör på \textbf{1 Mbit/s}, och bussen är en kabel på labbänken.
Gränsen är alltså aldrig i närheten av att bli ett problem --- men sambandet är värt att känna
till, eftersom det förklarar varför fordonssystem ofta kör långsammare än de skulle kunna.

\section*{Övningar}

\exercise{Kablage}{Räkning}
\label{ex:1:1}
\begin{enumerate}[a)]
  \item Hur många förbindelser krävs för att koppla ihop åtta enheter punkt-till-punkt?
  \item Hur många anslutningar krävs om de i stället delar en buss?
  \item Vid vilket antal enheter blir punkt-till-punkt dubbelt så många förbindelser som
    bussens anslutningar?
\end{enumerate}

\exercise{Bussens nivå}{Resonemang}
\label{ex:1:2}
Fyra noder sitter på en buss. Under ett visst bitintervall driver nod A recessivt, nod B
dominant, nod C recessivt, och nod D driver inte alls.
\begin{enumerate}[a)]
  \item Vilken nivå har bussen?
  \item Vilka av noderna läser tillbaka något annat än det de sände?
  \item Vad drar de för slutsats av det?
\end{enumerate}

\exercise{Terminering}{Resonemang}
\label{ex:1:3}
En grupp mäter 40~\textohm{} över en strömlös buss.
\begin{enumerate}[a)]
  \item Vad är fel?
  \item Bussen fungerar ändå på labbänken. Varför, och när skulle den sluta göra det?
\end{enumerate}

\exercise{Asymmetrin}{Resonemang}
\label{ex:1:4}
Anta en tänkt buss där recessiv i stället vann över dominant --- alltså där bussen är recessiv om
minst en nod driver den recessivt.
\begin{enumerate}[a)]
  \item Skulle arbitreringen fortfarande fungera?
  \item Vilket identifierarvärde skulle i så fall ha högst prioritet?
  \item Varför valde CAN den ordning den valde? Tänk på vad bussen gör när ingen sänder.
\end{enumerate}
